\documentclass[a4paper,12pt,UTF8]{ctexart}
\usepackage{tikz}
\usetikzlibrary{shapes.misc,arrows.meta,fit,positioning,backgrounds}
\usepackage[margin=0.5cm]{geometry}
\usepackage{amsmath}
\usepackage{amssymb}
\usepackage{fontspec}

\newfontfamily\mj[Scale=1.2]{Riichi-Mahjong}

\begin{document}

% \[
%     \left.\begin{aligned}
%         \text{老头牌：{\mj 1s 9s 1p 9p 1m 9m}}\\
%         \left.
%         \begin{aligned}
%             \text{风牌：{\mj 1z 2z 3z 4z}}\\
%             \text{三元牌：{\mj 5z 6z 7z}}
%         \end{aligned}\right\} \text{字牌}
%     \end{aligned}\right\} \text{幺九牌（幺）}
% \]

% 绿牌：{\mj 2s 3s 4s 6s 8s 6z}

% * 门清限定

\begin{tikzpicture}[scale=0.6, transform shape,
    declare function={x=3.5; y=-1.1;},
    yi/.style={
        rounded rectangle,
        draw,
        fill=white,
        minimum width=1.5cm,
        minimum height=0.8cm,
        font=\small
    },
    tag/.style={
        rounded rectangle,
        fill=gray!30,
        minimum width=1.5cm,
        minimum height=0.6cm,
        font=\small
    },
    arr/.style={
        ->,
        >=Stealth,
        font=\small,
        to path={|- (\tikztotarget) \tikztonodes}
    },
    groupbox/.style={
        rounded corners=5pt,
        inner sep=5pt,
        minimum width=2cm
    }
]

\node[tag] (t1) at (0, .3*y-y) {特殊顺子};
    \node[yi] (ssts) at (x, 0) {三色同顺};
    \node[yi] (*ssts) at (2*x, 0) {*三色同顺};
    \draw[arr] (t1) to node[near end, above] {三色同数字} (ssts);
    \draw[arr] (ssts) -- node[above] {门清} (*ssts);

    \node[yi] (yqtg) at (x, -y) {一气通贯};
    \node[yi] (*yqtg) at (2*x, -y) {*一气通贯};
    \node[yi] (jlbd) at (7*x, -y) {*九莲宝灯};
    \node[yi] (czjlbd) at (7*x, y) {*纯正九莲宝灯};
    \draw[arr] (t1) to node[near end, above] {从一到九} (yqtg);
    \draw[arr] (yqtg) -- node[above] {门清} (*yqtg);
    \draw[arr] (*yqtg) -- node[near start, above] {1和9各有3张} (jlbd);
    \draw[arr] (jlbd) -- node[right] {听9张} (czjlbd);

    \node[tag] (t2) at (0.4*x, y) {*杯口};
    \node[yi] (ybk) at (x, y) {*一杯口};
    \node[yi] (ebk) at (3*x, y) {*二杯口};
    \draw[arr] (t2) -- (ybk);
    \draw[arr] (t1) to node[near end, below] {拐子龙} (t2);
    \draw[arr] (ybk) -- node[above] {再来一对} (ebk);

\node[tag] (t3) at (0,2.3*y) {特殊花色};
    \node[yi] (lys) at (7*x, 2*y) {绿一色};
    \draw[arr] (t3) to node[near end, above] {全是绿牌} (lys);

    \node[yi] (hys) at (2*x, 3*y) {混一色};
    \node[yi] (*hys) at (3*x, 3*y) {*混一色};
    \node[yi] (qys) at (5*x, 4*y) {清一色};
    \node[yi] (*qys) at (6*x, 4*y) {*清一色};
    \draw[arr] (t3) to node[near end, above] {条饼万只占1种} (hys);
    \draw[arr] (hys) -- node[above] {门清} (*hys);
    \draw[arr] (*hys) to[to path={-| (\tikztotarget) \tikztonodes}] node[near start, above] {没有字牌} (*qys);
    \draw[arr] (hys) to node[near end, above] {没有字牌} (qys);
    \draw[arr] (qys) -- node[above] {门清} (*qys);

\node[tag] (t4-) at (0,4.3*y) {特殊刻杠};
\node[tag] (t4) at (0.5*x,5.3*y) {3组刻杠};
    \draw[arr] (t4-) to (t4);

    \node[yi] (ddh) at (2*x,5*y) {对对和};
    \node[circle, fill=black, minimum size=4pt, inner sep=0pt] (x1) at (6*x,5*y) {};
    \draw[arr] (t4) to node[near end, above] {再碰杠} (ddh);
    \draw (ddh) -- node[above] {门清} (x1);
    \node[yi] (4ak) at (7*x,5*y) {*四暗刻};
    \node[yi] (sakdq) at (7*x,6*y) {*四暗刻单骑};
    \draw[arr] (x1) -- node[above] {最后凑刻子} (4ak);
    \draw[arr] (x1) to node[near end, above] {最后凑雀头} (sakdq);

    \node[yi] (sstk) at (2*x,6*y) {三色同刻};
    \draw[arr] (t4) to node[near end, above] {三色同数字} (sstk);

    \node[yi] (3ak) at (2*x,7*y) {三暗刻};
    \draw[arr] (t4) to node[near end, above] {全是暗的} (3ak);

    \node[yi] (3gz) at (2*x,8*y) {三杠子};
    \node[yi] (4gz) at (7*x,8*y) {四杠子};
    \draw[arr] (t4) to node[near end, above] {全是杠} (3gz);
    \draw[arr] (3gz) -- node[above] {再杠} (4gz);

    \node[yi] (hlt) at (2*x,9*y) {混老头};
    \node[yi] (qlt) at (7*x,9*y) {清老头};
    \draw[arr] (t4) to node[near end, above] {全是幺} (hlt);
    \draw[arr] (hlt) -- node[above] {没有字牌} (qlt);

\node[tag] (t5) at (0,9.3*y) {幺};
\node[tag] (t6) at (0.5*x,10*y) {老头};
    \draw[arr] (t5) to (t6);
    \draw[arr] (t6) to[to path={-| (\tikztotarget) \tikztonodes}] node[near start, above] {七对幺} (hlt);

    \node[yi] (hqdyj) at (x, 11*y) {混全带幺九};
    \draw[arr] (t5) to node[near end, above] {全带幺} (hqdyj);
    \draw[arr] (hqdyj) to[to path={-| (\tikztotarget) \tikztonodes}] node[near start, above] {全是刻杠} (hlt);
    \node[yi] (*hqdyj) at (2*x, 12*y) {*混全带幺九};
    \draw[arr] (hqdyj) to node[near end, above] {门清} (*hqdyj);
    \node[yi] (cqdyj) at (2*x, 13*y) {纯全带幺九};
    \draw[arr] (hqdyj) to node[near end, above] {没有字牌} (cqdyj);
    \node[yi] (*cqdyj) at (3*x, 13*y) {*纯全带幺九};
    \draw[arr] (cqdyj) -- node[above] {门清} (*cqdyj);
    \draw[arr] (*hqdyj) to[to path={-| (\tikztotarget) \tikztonodes}] node[near start, above] {没有字牌} (*cqdyj);

    \node[tag] (t7) at (.5*x, 14*y) {三元};
    \node[yi] (xsy) at (2*x, 14*y) {小三元};
    \node[yi] (dsy) at (7*x, 14*y) {大三元};
    \draw[arr] (t7) -- node[above] {{\mj 5z5z5z6z6z6z7z7z7z}凑8张} (xsy);
    \draw[arr] (xsy) -- node[above] {凑齐} (dsy);
    \draw[arr] (t5) to (t7);

    \node[tag] (t8) at (.5*x, 15*y) {四喜（风）};
    \node[yi] (xsx) at (7*x, 15*y) {小四喜};
    \node[yi] (dsx) at (7*x, 17*y) {大四喜};
    \draw[arr] (t8) -- node[above] {{\mj 1z1z1z2z2z2z3z3z3z4z4z4z}凑11张} (xsx);
    \draw[arr] (xsx) -- node[right] {凑齐} (dsx);
    \draw[arr] (t5) to (t8);

    \node[yi] (dyj) at (x, 16*y) {断幺九};
    \draw[arr] (t5) to node[near end, above] {没有幺} (dyj);

    \node[yi] (y) at (x, 17*y) {役};
    \draw[arr] (t5) to node[near end, above] {自风/场风/三元} (y);

    \node[yi] (ljmg) at (4*x, 18*y) {流局满贯};
    \draw[arr] (t5) to node[near end, above] {只打幺，不被吃碰杠，直到流局} (ljmg);

\node[tag] (t9) at (0, 19*y) {*其他门清限定};
    \node[yi] (ph) at (x, 19*y) {*平和};
    \node[yi] (lz) at (x, 20*y) {*立直};
    \node[yi] (qdz) at (2*x, 19*y) {*七对子};

\node[tag] (t10) at (0, 22*y) {看运气};
    \node[yi] (yf) at (x, 22*y) {*一发};
    \draw[arr] (lz) -- node[right] {1巡内} (yf);
    \node[yi] (qg) at (x, 24*y) {抢杠};
    \node[yi] (mqqzmh) at (x, 23*y) {*门前清自摸和};
    \node[tag] (t11) at (.3*x, 25.3*y) {最后1张};
    \node[yi] (hdly) at (x, 25*y) {海底捞月};
    \node[yi] (hdmy) at (x, 26*y) {河底摸鱼};
    \draw[arr] (t11) to node[near end, above] {自摸} (hdly);
    \draw[arr] (t11) to node[near end, below] {荣和} (hdmy);

    \node[yi] (gsws) at (7*x, 22*y) {*国士无双};
    \node[yi] (gswsssm) at (7*x, 24*y) {*国士无双十三面};
    \draw[arr] (gsws) -- node[right] {听13张} (gswsssm);
    \node[yi] (th) at (7*x, 25*y) {*天和};
    \node[yi] (dh) at (7*x, 26*y) {*地和};


\node[font=\small\bfseries] (1f) at (x, -2*y) {1番} ;
\node[font=\small\bfseries] (2f) at (2*x, -2*y) {2番} ;
\node[font=\small\bfseries] (3f) at (3*x, -2*y) {3番} ;
\node[font=\small\bfseries] (4f) at (4*x, -2*y) {满贯} ;
\node[font=\small\bfseries] (5f) at (5*x, -2*y) {5番} ;
\node[font=\small\bfseries] (6f) at (6*x, -2*y) {6番} ;
\node[font=\small\bfseries] (7f) at (7*x, -2*y) {役满} ;

\begin{pgfonlayer}{background}
% 6字及以上的胶囊没有纳入，因为会撑宽方块，影响美观。
\node[groupbox, fit=(yqtg)(ssts)(ybk)(hqdyj)(dyj)(y)(ph)(lz)(yf)(qg)(hdly)(hdmy)(1f), fill=red!10] {};
\node[groupbox, fit=(*yqtg)(*ssts)(hys)(ddh)(sstk)(3ak)(3gz)(hlt)(*hqdyj)(cqdyj)(xsy)(qdz)(2f), fill=orange!10] {};
\node[groupbox, fit=(ebk)(*hys)(*cqdyj)(3f), fill=yellow!10] {};
\node[groupbox, fit=(ljmg)(4f), fill=green!10] {};
\node[groupbox, fit=(qys)(5f), fill=cyan!10] {};
\node[groupbox, fit=(*qys)(6f), fill=blue!10] {};
\node[groupbox, fit=(jlbd)(lys)(4ak)(sakdq)(4gz)(qlt)(dsy)(xsx)(dsx)(gsws)(th)(dh)(7f), fill=purple!10] {};
\end{pgfonlayer}

\end{tikzpicture}

\end{document}
